\section*{Chapitre \ref{ch:b1:clastic-sediments} --- Les sédiments clastiques et leurs milieux}

\begin{solution}{exo:b1:clastic-sediments:1}
$\varphi = -\log_2 d$~: \qty{0.18}{mm} donne $2.47$, un sable fin (2 à 3)~;
\qty{0.75}{mm} donne $0.42$, un sable grossier (0 à 1)~; \qty{3}{mm} donne
$-1.58$, un granule ($-1$ à $-2$)~; \qty{0.012}{mm} donne $6.38$, un silt fin (6
à 7)~; \qty{150}{mm} donne $-7.23$, un galet ($-6$ à $-8$).
\end{solution}

\begin{solution}{exo:b1:clastic-sediments:2}
$M_z = (1.05 + 1.71 + 2.43)/3 = 1.73\,\varphi$, c'est-à-dire $2^{-1.73} =
\qty{0.30}{mm}$, un sable moyen. $\sigma_I = (2.43 - 1.05)/4 + (2.98 - 0.62)/6.6
= 0.345 + 0.358 = 0.70\,\varphi$~: modérément bien classé (0.50 à 0.71), juste
sous la limite du modérément classé.
\end{solution}

\begin{solution}{exo:b1:clastic-sediments:3}
$q_b = (1 - 0.4) \times 0.02 \times (\num{0.01}/60)/2 = \qty{1.0e-6}{m^2.s^{-1}}$
de grains solides. En un jour, $\num{1.0e-6} \times 86400 = \qty{0.0864}{m^3}$ par
mètre de crête, soit $0.0864 \times 2650 = \qty{229}{kg}$ de quartz.
\end{solution}

\begin{solution}{exo:b1:clastic-sediments:4}
(a)~50 nœuds font $50 \times 1852/3600 = \qty{25.7}{m.s^{-1}}$~; 12 nœuds,
\qty{6.17}{m.s^{-1}}. (b)~$\qty{100000}{m}/\qty{6.17}{m.s^{-1}} =
\qty{16200}{s}$, environ \qty{4.5}{h}. (c)~480 milles font \qty{772}{km}, 400 milles
\qty{644}{km}.
\end{solution}

\begin{solution}{exo:b1:clastic-sediments:5}
(a)~$(\varphi_{84} - \varphi_{16})/2 = 2 \times 0.9945\,\sigma/2 =
0.9945\,\sigma$, and $(\varphi_{95} - \varphi_5)/3.3 = 2 \times
1.6449\,\sigma/3.3 = 0.9969\,\sigma$~: les deux estiment $\sigma$ à un tiers
de pour cent près. (b)~Un sédiment réel n'a pas à être normal en unités phi~: un petit
apport de grains bien plus grossiers ou bien plus fins change beaucoup $\varphi_5$ ou
$\varphi_{95}$ et peu $\varphi_{16}$ et $\varphi_{84}$, donc la seconde
estimation croît alors que la première non. Les deux diffèrent le plus pour les mélanges, tels
qu'un sable avec un peu de gravier et de boue, ou un dépôt bimodal.
\end{solution}

\begin{solution}{exo:b1:clastic-sediments:6}
(a)~$\sin\theta = T/L = 1.5/15 = 0.1$, donc $\theta = \ang{5.74}$. (b)~$a =
c\tan\theta = 0.5 \times 0.1005 = \qty{0.050}{cm}$ par minute, soit
\qty{3.0}{cm} par heure. (c)~Un fond exhaussé de centimètres par heure reçoit du sable
tombant de la suspension très vite~: un courant très chargé qui ralentit,
comme une crue qui décroît ou la queue d'un \omterm{def:b1:clastic-sediments:turbidite}{courant de turbidité}.
\end{solution}

\begin{solution}{exo:b1:clastic-sediments:7}
(a)~$(350 + 20 + 5 + 340 + 15 + 30)/6 = 760/6 = 126.7$~: une direction
sud-est, bien que toutes les lectures soient à moins de \ang{20} du nord. (b)~$S =
\sum\sin\alpha_i = 0.672$, $C = \sum\cos\alpha_i = 5.692$, so $\bar\alpha =
\operatorname{atan2}(0.672, 5.692) = \ang{6.7}$, and $\bar R =
\sqrt{0.672^2 + 5.692^2}/6 = 5.732/6 = 0.955$. (c)~La moyenne arithmétique traite
350 et 005 comme éloignés, alors que sur le cercle ils sont à \ang{15}~: les \omterm{def:b1:spherical-astronomy:horizontal}{azimuts}
doivent être moyennés comme des vecteurs unitaires. Le grand $\bar R$ montre que les lectures sont
étroitement groupées autour de \ang{7}.
\end{solution}

\begin{solution}{exo:b1:clastic-sediments:8}
A a plus de \qty{5}{\percent} d'argile~: immature, et immature aussi du point de vue de la composition,
avec son \omterm{def:b1:rock-forming-minerals:groups}{feldspath} et ses fragments de roche~; déposé rapidement près de sa
source, sur un cône alluvial ou une plaine d'inondation, ou par un \omterm{def:b1:clastic-sediments:turbidite}{courant de turbidité}. B a moins
de \qty{5}{\percent} d'argile et est bien classé ($\sigma_I < 0.5\,\varphi$) mais ses
grains sont subanguleux ($\rho < 3$)~: mature, par exemple une barre de rivière ou un sable
de plateau vanné par les courants. C n'a pas d'argile, est très bien classé et ses grains sont
arrondis ($\rho > 3$)~: supermature, et avec \qty{98}{\percent} de quartz aussi
mature du point de vue de la composition, comme sur une plage ou dans un champ de \omterm{def:b1:clastic-sediments:bedform}{dunes} du désert.
\end{solution}

\begin{solution}{exo:b1:clastic-sediments:9}
(a)~$1.5/\sqrt{9.81 \times 0.3} = 0.87$~; $1.0/\sqrt{9.81 \times 2} = 0.23$.
(b)~$h = U^2/g = 1.5^2/9.81 = \qty{0.23}{m}$~; $U = \sqrt{9.81 \times 2} =
\qty{4.43}{m.s^{-1}}$. (c)~Le nombre de Froude croît quand la profondeur diminue à
vitesse donnée~: de minces nappes d'eau rapides, comme le retrait d'une vague sur une plage,
atteignent facilement des nombres de Froude proches de 1, alors qu'une rivière profonde de plusieurs mètres
demanderait plusieurs mètres par seconde, que peu de rivières atteignent sauf en crue.
\end{solution}

\begin{solution}{exo:b1:clastic-sediments:10}
(a)~Avec les milieux $-0.5, 0.5, \ldots, 4.5$~: $m = (-1 + 4 + 45 + 100 + 52.5
+ 22.5)/100 = 2.23\,\varphi$~; les carrés des écarts pondérés par les
pourcentages font en tout \num{107.71}, donc $s = \sqrt{1.0771} = 1.04\,\varphi$~:
mal classé. (b)~Le pourcentage cumulé plus grossier est 0, 2, 10, 40, 80, 95 et
100 à $-1, 0, \ldots, 5\,\varphi$. Alors $\varphi_5 = 0 + 3/8 = 0.375$,
$\varphi_{16} = 1 + 6/30 = 1.20$, $\varphi_{50} = 2 + 10/40 = 2.25$,
$\varphi_{84} = 3 + 4/15 = 3.27$, $\varphi_{95} = 4.0$. $M_z = (1.20 + 2.25 +
3.27)/3 = 2.24\,\varphi$ et $\sigma_I = (3.27 - 1.20)/4 + (4.0 - 0.375)/6.6 =
0.517 + 0.549 = 1.07\,\varphi$. (c)~Les deux méthodes s'accordent à environ
\num{0.01} sur la moyenne et \num{0.03} sur le \omterm{def:b1:clastic-sediments:sorting}{classement}~: un sable fin
($2.24\,\varphi$ fait \qty{0.21}{mm}), mal classé.
\end{solution}

\begin{solution}{exo:b1:clastic-sediments:11}
(a)~$U = \sqrt{0.02 \times 9.81 \times 50 \times 0.005/0.003} =
\qty{4.04}{m.s^{-1}}$, soit 7.9 nœuds. (b)~$(\Delta\rho/\rho)\,h = U^2
C_f/(gS) = 25.7^2 \times 0.003/(9.81 \times 0.005) = \qty{40}{m}$~: par exemple
un courant épais de \qty{400}{m} avec un excès de densité de \qty{10}{\percent}, ce qui
montre combien le courant de 1929 devait être grand et dense. (c)~La vitesse
varie comme $\sqrt S$~: $4.04/\sqrt 5 = \qty{1.81}{m.s^{-1}}$. Le courant plus lent
a une turbulence plus faible et laisse ses grains les plus grossiers~; il devient moins dense, ralentit
encore, et étale sa charge sur la plaine en un \omterm{def:b1:clastic-sediments:turbidite}{banc granoclassé}.
\end{solution}

\begin{solution}{exo:b1:clastic-sediments:12}
(a)~$0.7 \times \qty{5e10}{kg}/\qty{2000}{kg.m^{-3}} = \qty{1.75e7}{m^3}$ par
an. (b)~Étalé sur un front haut de \qty{30}{m} et large de \qty{100}{km}, il
fait avancer le front de $\num{1.75e7}/(30 \times \num{1e5}) = \qty{5.8}{m}$ par
an. (c)~$a = v\tan\alpha = 5.83 \times \tan\ang{0.5} = \qty{0.051}{m}$, environ
\qty{5}{cm} par an, cent fois moins que l'avancée~; la
succession de \qty{30}{m} se forme en $30/0.051 = 590$ ans.
\end{solution}

\begin{solution}{pb:b1:clastic-sediments:1}
\textbf{1.} Unité~1~: $(5.0 + 6.6 + 8.4)/3 = 6.67\,\varphi$, \qty{0.0098}{mm}
(\qty{9.8}{\micro m})~; unité~2~: $3.47\,\varphi$, \qty{0.090}{mm}~; unité~3~:
$1.80\,\varphi$, \qty{0.29}{mm}~; unité~4~: $0.73\,\varphi$, \qty{0.60}{mm}.

\textbf{2.} Unité~1~: $(8.4 - 5.0)/4 + (9.6 - 4.2)/6.6 = 0.850 + 0.818 =
1.67\,\varphi$~; unité~2~: $0.550 + 0.545 = 1.10$~; unité~3~: $0.275 + 0.273 = 0.55$~;
unité~4~: $0.725 + 0.712 = 1.44$.

\textbf{3.} Unité~1~: silt fin, mal classé~; unité~2~: sable très fin, mal
classé~; unité~3~: sable moyen, modérément bien classé~; unité~4~: sable grossier,
mal classé.

\textbf{4.} Unité~3. À l'embouchure du fleuve le courant et les vagues déplacent le sable
sans cesse et emportent les fines, tandis que la boue du prodelta se dépose depuis la
suspension en un mélange de toutes les tailles fines qui y parviennent.

\textbf{5.} La taille $8\,\varphi$ se trouve aux $1.4/1.8$ du chemin de
$\varphi_{50}$ à $\varphi_{84}$, donc $50 + (1.4/1.8) \times 34 = 76$ pour cent de
l'échantillon est plus grossier et environ \qty{24}{\percent} est de la taille de l'argile. De telles particules
fines ne se déposent que depuis une eau presque calme~: le site était sous la portée
des vagues et des courants.

\textbf{6.} Un chenal transporte à la fois du gravier, roulé sur son fond lors des crues, et du
sable~; le dépôt mêle le dépôt résiduel de gravier de l'écoulement le plus profond au sable des
écoulements plus faibles, donc il couvre de nombreuses classes.

\textbf{7.} La moyenne passe de \qty{0.0098}{mm} à \qty{0.090}{mm} et
\qty{0.29}{mm}~: la succession se grossit vers le haut.

\textbf{8.} Unité~1, boue laminée~: le prodelta~; unité~2, alternance de silt et de
sable~: le bas du front de delta~; unité~3, sable à \omterm{def:b1:clastic-sediments:cross-bedding}{stratification entrecroisée}~: les barres d'embouchure du
haut du front de delta.

\textbf{9.} Chaque zone est venue reposer sur celle qui était auparavant devant elle~:
le delta s'est construit vers l'avant, vers le large, sur le site.

\textbf{10.} Au moins \qty{21}{m}, l'épaisseur des unités 1 à 3, puisque la
succession a rempli l'eau du fond du prodelta à la barre d'embouchure
près de la surface~; davantage, puisque les boues se sont compactées depuis
(\cref{ch:b1:diagenesis-basins}).

\textbf{11.} Un chenal distributaire de la plaine deltaïque, entaillant la barre
d'embouchure à mesure que le delta se construisait vers l'avant~; ses cailloux sont sa \omterm{def:b1:erosion-transport:load}{charge de fond}.

\textbf{12.} Non~: une barre de méandre commence par une surface d'érosion et un dépôt résiduel et s'affine
vers le haut, jusqu'à la boue de plaine d'inondation. Le front de delta se grossit vers le haut à partir de la boue.

\textbf{13.} $\qty{16}{m}/\num{4000}\text{ ans} = \qty{4.0}{mm}$ par an.

\textbf{14.} $v = a/\tan\alpha = 0.0040/\tan\ang{0.5} = 0.0040/0.008727 =
\qty{0.458}{m}$ par an.

\textbf{15.} $0.458 \times 4000 = \qty{1830}{m}$, environ \qty{1.8}{km}.

\textbf{16.} $\qty{21}{m}/\qty{4.0}{mm.\year^{-1}} = 5250$ ans.

\textbf{17.} Le front de delta reçoit l'essentiel du sable près des embouchures, le
prodelta bien moins et la plaine deltaïque peu sauf en crue~; les embouchures actives
se déplacent aussi d'un lieu à l'autre, si bien qu'un point peut recevoir beaucoup pendant un temps puis
presque rien. Un taux unique entre deux bancs datés est une moyenne sur ces
changements, et les boues ont perdu une partie de leur épaisseur par compaction.

\textbf{18.} La somme est 1544, donc la moyenne arithmétique est \ang{128.7}, une
direction sud-est~: dénuée de sens, puisque les lectures se dispersent autour du nord
et que les valeurs proches de 340 et de 010 sont moyennées comme si elles étaient éloignées.

\textbf{19.} $S = \sum\sin\alpha_i = 1.698$ et $C = \sum\cos\alpha_i =
11.248$.

\textbf{20.} $\bar R = \sqrt{1.698^2 + 11.248^2}/12 = 11.376/12 = 0.948$~: les
lectures sont étroitement groupées autour d'une seule direction.

\textbf{21.} Les rivières coulaient, et le delta se construisait, vers environ le nord
(\cref{prop:b1:clastic-sediments:azimuths})~; la terre se trouvait au sud.

\textbf{22.} Parce que les feuillets frontaux ont été mesurés sur des bancs qui ne sont plus
horizontaux~: leurs sens de \omterm{def:b1:geological-maps:attitude}{pendage} sont ceux des bancs inclinés plus
l'inclinaison. Faire tourner chaque feuillet frontal autour de la direction de la stratification de l'angle de \omterm{def:b1:geological-maps:attitude}{pendage} de
la stratification restaure la stratification, et avec elle les feuillets frontaux, dans leur
\omterm{def:b1:geological-maps:attitude}{attitude} d'origine.

\textbf{23.} Les courants de marée vont et viennent, et construisent des ensembles qui plongent en deux
sens opposés, comme dans la \omterm{def:b1:clastic-sediments:cross-bedding}{stratification entrecroisée} en arêtes de poisson~; le diagramme en rose aurait
deux lobes à environ \ang{180} l'un de l'autre. Les douze lectures forment un seul lobe,
$\bar R$ proche de 1~: un courant de rivière s'écoulant dans un seul sens.

\textbf{24.} $\bar\alpha = \operatorname{atan2}(1.698, 11.248) = \ang{8.6}$~:
l'\omterm{def:b1:spherical-astronomy:horizontal}{azimut} moyen du \omterm{def:b1:clastic-sediments:cross-bedding}{paléocourant} de l'unité~3 est \ang{9}, un peu à l'est du nord.
\end{solution}
